% Appendice: Repertorio dei siti preistorici veneti nel catalogo MiC
% Abbreviazione usata in tutto il capitolo:
%   a.p.z.a. = anni prima dell'anno zero astronomico

\chapter{Repertorio dei siti preistorici veneti nel catalogo MiC}
\label{cap:appendice_siti}

\aprecapitolo{I}

\section*{Nota introduttiva}

I dati raccolti in questa appendice provengono interamente dal catalogo ufficiale del Ministero della Cultura italiano, accessibile al pubblico attraverso il portale \textit{catalogo.beniculturali.it} gestito dall'Istituto Centrale per il Catalogo e la Documentazione (ICCD). Le schede sono state consultate e rielaborate nell'estate del 2026 ed erano quelle disponibili a quella data.

Il catalogo MiC non è un inventario sistematico di tutto ciò che esiste nel sottosuolo veneto: è il registro di ciò che è stato formalmente catalogato, cioè scavato, segnalato e inserito nel sistema nazionale da parte di soprintendenze, università, musei civici e privati cittadini con le debite autorizzazioni. La sua incompletezza è dunque strutturale, ma anche istruttiva: le zone con densità alta di schede corrispondono ad aree esplorate; le zone vuote ci dicono dove la ricerca non ha ancora guardato.

Abbiamo incluso tutti i 136 siti con almeno una fase databile a età preistorica (dal Paleolitico all'età del Ferro), escludendo i siti esclusivamente romani o medievali (una trentina di schede). I siti sono ordinati per epoca, dalla più antica alla più recente, e all'interno di ciascuna epoca per zona geografica. Per i siti che attraversano più epoche, la descrizione completa si trova alla prima occorrenza; le sezioni successive riportano solo un rimando. Per le date numeriche usiamo l'abbreviazione \textbf{a.p.z.a.} (anni prima dell'anno zero astronomico), equivalente a ciò che la letteratura scientifica indica comunemente con ``a.C.'', ma senza l'ambiguità religiosa di quella formula.

% ============================================================
\section{Paleolitico}
\label{app:paleolitico}
% 13 siti
% ============================================================

Il Paleolitico copre un intervallo enorme, dai primi manufatti musteriani attribuibili all'uomo di Neanderthal (circa 300.000 a.p.z.a.) fino all'Epigravettiano finale (circa 10.000 a.p.z.a.), quando i cacciatori-raccoglitori del Paleolitico superiore si adattano al progressivo scioglimento dei ghiacciai. Nel Veneto questo periodo è documentato principalmente nella Lessinia veronese e, in misura minore, nell'Altopiano di Asiago e in alcune grotte dei Colli Berici.

\subsection{Altopiano Vicentino}

\textbf{Roana (VI).} Una grotta nel territorio di Roana ha restituito un deposito litico musteriano, senza ulteriori specificazioni nel catalogo. La scheda è essenziale: il sito è noto ma non scavato sistematicamente.

\textbf{Asiago (VI).} Un riparo sotto roccia in territorio di Asiago ha fornito evidenze di frequentazione epigravettiana finale, interpretate come sosta stagionale di cacciatori nell'area altopianare, probabilmente legata ai movimenti della fauna in quota.

\subsection{Colli Euganei e Berici}

\textbf{Longare (VI).} La grotta nel territorio di Longare mostra una sequenza che abbraccia Musteriano, Aurignaziano e Gravettiano, suggerendo frequentazioni ripetute in un arco temporale molto lungo. L'interpretazione MiC la inquadra come rifugio stagionale; si tratta di uno dei pochi siti dei Berici con attestazioni paleolitiche multiple.

\textbf{Cinto Euganeo (PD).} Nei pressi di Valnogaredo, i ritrovamenti superficiali sono interpretati come tracce di frequentazione musteriana. Non sono state individuate strutture; la localizzazione ai margini dei Colli Euganei è geograficamente interessante per comprendere le rotte di caccia in area collinare.

\subsection{Lessinia}

\textbf{Dolcè (VR).} Un riparo sotto roccia nel fondovalle dell'Adige, scavato tra il 1983 e il 1988, conserva una sequenza che va dal Paleolitico superiore al Mesolitico (culture Sauveterriana e Castelnoviana). Nel 1983 fu segnalata anche una sepoltura a inumazione di probabile età olocenica. Il sito è quindi multifase: vedi anche la sezione Mesolitico~(\ref{app:mesolitico}).

\textbf{Fumane (VR) -- Grotta di Fumane.} \label{app:fumane_paleo} Uno dei siti paleolitici più importanti d'Europa, secondo la valutazione stessa del catalogo MiC. La grotta, alla base della Valle dei Progni, fu intravista nel tardo Ottocento durante lavori stradali e saccheggiata ripetutamente prima che, a partire dal 1988, un gruppo di ricerca articolato tra Soprintendenza, Università di Ferrara, Università di Milano e Museo di Scienze Naturali di Verona avviasse scavi sistematici. La sequenza stratigrafica comprende livelli riferibili al Paleolitico medio (Musteriano), al Paleolitico superiore iniziale (Uluzziano, poi Aurignaziano) e al Gravettiano. La presenza di manufatti su osso, ornamenti personali e possibili tracce di comportamento simbolico colloca questa grotta tra i contesti di riferimento per la transizione tra Neanderthal e uomo anatomicamente moderno in Italia nord-orientale. Per la fase neolitica ed eneolitica vedi anche la scheda nella sezione Neolitico~(\ref{app:fumane_neo}).

\textbf{Grezzana (VR).} Un sito classificato come ``uno dei più importanti d'Italia'' per ampiezza stratigrafica, quantità di reperti e qualità della documentazione, con sequenze che vanno dal Paleolitico medio musteriano all'Aurignaziano e all'Epigravettiano. Non si hanno indicazioni di apertura turistica o di scavi recenti nel catalogo.

\textbf{Sant'Anna d'Alfaedo (VR).} Scavi condotti a partire dal 1922, sia nella caverna che nelle aree circostanti, hanno documentato una continuità di presenza umana dal Paleolitico inferiore al Neolitico. Il catalogo segnala come particolarmente sviluppata la lavorazione della selce, attiva sino all'età del rame. Vedi la sezione Neolitico~(\ref{app:alfaedo_neo}).

\textbf{Sant'Ambrogio di Valpolicella (VR).} Tracce di frequentazione musteriana, senza ulteriori dettagli nella scheda. La posizione all'imbocco della Valpolicella suggerisce un possibile punto di caccia o di transito lungo il bordo meridionale della Lessinia.

\subsection{Pianura Veronese e Bassa}

\textbf{Roncà (VR).} \label{app:ronca_paleo} Sito pluristratificato con fase musteriana, seguita da fasi eneolitica, bronzo recente-finale e tarda età del ferro, con tracce di romanizzazione. La scheda è sintetica (``sito pluristratificato'') senza descrizione delle strutture. Vedi anche le sezioni Bronzo~(\ref{app:ronca_bronzo}) e Ferro~(\ref{app:ronca_ferro}).

\subsection{Pianura e Delta}

\textbf{Vo (PD).} Frequentazione musteriana individuata nel Padovano. Tracce di superficie senza strutture accertate; la scheda interpreta il rinvenimento come frequentazione stagionale dell'area da parte di gruppi neanderthaliani.

\subsection{Prealpi Trevigiane e Pedemontana}

\textbf{Tarzo (TV).} \label{app:tarzo_paleo} Sito pluristratificato con tracce dell'uomo di Neanderthal (Paleolitico medio), frequentazione mesolitica castelnoviana e tracce di età del Bronzo. Il sito si trova nel versante meridionale delle Prealpi orientali; la successione di popolamenti ne fa un punto di riferimento per la preistoria trevigiana. Vedi anche le sezioni Mesolitico~(\ref{app:mesolitico}) e Bronzo~(\ref{app:tarzo_bronzo}).

\subsection*{Nota di sintesi -- Paleolitico}

I tredici siti paleolitici veneti conosciuti sono distribuiti in modo tutt'altro che casuale: la Lessinia veronese concentra i siti più ricchi e meglio documentati, grazie alla disponibilità di grotte e ripari in una roccia calcarea abbondante di selce di buona qualità. La pianura padana e il delta sono quasi assenti, non perché fossero disabitate ma perché migliaia di anni di depositi alluvionali ricoprono i livelli paleolitici a profondità difficilmente raggiungibili dagli scavi ordinari. Le Prealpi trevigiane e l'Altopiano di Asiago mostrano frequentazioni marginali, probabilmente legate a percorsi di caccia stagionale verso le quote medie. Il Bellunese e il Cadore sono del tutto assenti da questa lista, fatto che riflette la scarsità di prospezioni sistematiche in quota, non necessariamente l'assenza di occupazione umana.

% ============================================================
\section{Mesolitico}
\label{app:mesolitico}
% 3 siti
% ============================================================

Il Mesolitico (grossomodo tra 10.000 e 6.000 a.p.z.a.) è il periodo di adattamento alle condizioni postglaciali: i boschi avanzano, la fauna cambia, le strategie di sussistenza si moltiplicano. Nel Veneto il Mesolitico è documentato quasi esclusivamente in zone di transizione tra pianura e montagna, lungo fiumi e laghi.

\subsection{Lessinia}

\textbf{Dolcè (VR).} Vedi la descrizione completa alla sezione Paleolitico~(\ref{app:paleolitico}). La fase mesolitica comprende culture Sauveterriana e Castelnoviana, indicando una continuità di occupazione del riparo dal Paleolitico superiore attraverso l'intero Mesolitico.

\subsection{Prealpi Trevigiane e Pedemontana}

\textbf{Tarzo (TV).} Vedi la descrizione completa alla sezione Paleolitico~(\ref{app:tarzo_paleo}). La componente mesolitica è riferibile alla cultura Castelnoviana, ben attestata nel versante meridionale delle Prealpi orientali.

\subsection{Prealpi Vicentine e Piccole Dolomiti}

\textbf{Valdagno (VI).} Un riparo sotto roccia nel territorio di Valdagno ha restituito evidenze di frequentazione mesolitica. I saggi di scavo non hanno individuato strutture permanenti; si ipotizza un utilizzo stagionale del riparo, probabilmente connesso a caccia o pascolo nelle valli prealpine vicentine.

\subsection*{Nota di sintesi -- Mesolitico}

Tre sole schede per il Mesolitico veneto sono un numero che riflette soprattutto la storia della ricerca: il Mesolitico è difficile da identificare in superficie, i suoi siti sono spesso piccoli e con scarsa densità di materiali, e le prospezioni sistematiche in quota sono rare. I siti noti si concentrano lungo il bordo prealpino e nelle valli di accesso alla montagna, coerentemente con le strategie di mobilità stagionale dei cacciatori mesolitici. La rete di siti mesolitici documentata nel Trentino e nell'Alto Adige contigui (Riparo Dalmeri, Mondeval de Sora, piani di Eraclea) suggerisce che anche il Veneto montano fosse frequentato durante questo periodo, ma le tracce non sono ancora state trovate o catalogate.

% ============================================================
\section{Neolitico}
\label{app:neolitico}
% 22 siti
% ============================================================

Il Neolitico veneto (circa 6.000-3.500 a.p.z.a.) è caratterizzato dall'arrivo dell'agricoltura e dell'allevamento, con comunità che si stabilizzano in insediamenti più duraturi. La cultura dei ``Vasi a Bocca Quadrata'' è il marcatore culturale dominante nel Neolitico medio-recente della pianura padano-veneta.

\subsection{Altopiano Vicentino}

\textbf{Lusiana (VI).} Insediamento d'altura con elementi strutturali legati a un sistema di fortificazione, classificato come ``insediamento fortificato'' nell'ambito proto-veneto neolitico. La posizione elevata suggerisce una funzione di controllo del territorio, forse di frontiera tra aree di influenza culturale diversa.

\subsection{Bellunese e Feltrino}

\textbf{Fonzaso (BL) -- Malga Campon di Monte Avena.} \label{app:avena_fonzaso} Il sito di Malga Campon, da considerarsi parte di un unico contesto con il sito di Pedavena (vedi sotto), si estende sulla cima del Monte Avena a circa 1.400~m di quota. Le evidenze più antiche risalgono al Paleolitico (Musteriano e Aurignaziano), legate allo sfruttamento della selce locale e alla caccia. La fase eneolitica, con una probabile prosecuzione dello sfruttamento della selce, è seguita da un'occupazione medievale legata all'alpeggio. Il catalogo classifica il sito anche come neolitico, verosimilmente per la continuità con la fase eneolitica. Vedi anche la sezione Eneolitico~(\ref{app:eneolitico}).

\textbf{Pedavena (BL) -- Casera Campet di Monte Avena.} Estensione orientale del sito di Fonzaso nello stesso comune limitrofo. Le caratteristiche e la cronologia sono identiche (vedi la scheda di Malga Campon, \ref{app:avena_fonzaso}): la separazione nel catalogo riflette solo il confine comunale, non una distinzione funzionale o cronologica.

\subsection{Cadore e Comelico}

\textbf{Selva di Cadore (BL) -- Riparo Mandriz.} \label{app:mandriz} Un riparo sotto roccia frequentato come ricovero stagionale tra il Neolitico tardo e l'Eneolitico da comunità dedite alla pastorizia transumante e alla caccia. La quota e il contesto montano indicano che anche in quest'epoca le valli cadorine erano percorse da gruppi mobili che sfruttavano i pascoli estivi. Vedi anche la sezione Eneolitico~(\ref{app:eneolitico}).

\subsection{Colli Euganei e Berici}

\textbf{Arcugnano (VI) -- Lago di Fimon (tre schede).} \label{app:fimon} Il lago di Fimon, nelle colline Beriche, ha restituito tracce di almeno tre settori distinti di un grande abitato perilacustre neolitico (cultura dei Vasi a Bocca Quadrata, ambito proto-veneto). Le tre schede separate nel catalogo corrispondono a settori con funzioni diverse: abitativa, artigianale e funeraria. Il complesso nel suo insieme costituisce uno dei più importanti insediamenti neolitici dell'Italia nord-orientale, con una notevole durata di occupazione e una struttura spaziale differenziata.

\textbf{Longare (VI).} Materiali riconducibili al Neolitico (ambito proto-veneto) raccolti in superficie; in assenza di strutture individuate, si interpreta come frequentazione a carattere prevalentemente abitativo. Il sito si aggiunge al quadro della presenza neolitica nei Berici.

\textbf{Pojana Maggiore (VI).} Insediamento neolitico in pianura, attribuibile a un abitato sulla base della tipologia del materiale rinvenuto. La scheda è sintetica.

\textbf{Este (PD).} \label{app:este} Sito pluristratificato emerso presso lo scolo delle Monache, con fasi dal Neolitico finale all'età del Bronzo. L'insediamento aveva una zona in ambiente umido (probabilmente su palafitte) e una in asciutto su dosso. La scansione cronologica è difficile da precisare, ma l'occupazione inizia con una fase neo-eneolitica seguita da almeno tre fasi del Bronzo con diversa organizzazione spaziale. Vedi anche le sezioni Eneolitico~(\ref{app:eneolitico}) e Bronzo~(\ref{app:este_bronzo}).

\subsection{Lago di Garda e Adige}

\textbf{Garda (VR) -- Monte Luppia e dintorni.} \label{app:garda_neo} Sito pluristratificato di straordinaria continuità che abbraccia il Neolitico (cultura dei Vasi a Bocca Quadrata), l'Eneolitico, l'età del Ferro (Veneti antichi e seconda età del Ferro), l'età romano-tarda, il periodo altomedievale con sepolture longobarde, il Medioevo scaligero e l'epoca contemporanea con gallerie e trincee della prima guerra mondiale. La scheda non specifica la natura delle evidenze neolitiche ed eneolitiche, ma la stratificazione eccezionale del promontorio gardesano è ben nota dalla letteratura regionale. Per la fase più importante ai fini di questo libro, quella protostorica con i petroglifi, vedi la sezione Bronzo~(\ref{app:garda_petroglifi}). L'occupazione registrata si estende fino al periodo contemporaneo.

\subsection{Lessinia}

\textbf{Fumane (VR) -- Seconda scheda.} \label{app:fumane_neo} Seconda scheda MiC per Fumane, relativa a un sito diverso dalla Grotta di Fumane paleolitica (vedi \ref{app:fumane_paleo}): insediamento e area funeraria con fasi neolitiche (cultura dei Vasi a Bocca Quadrata) ed eneolitiche, scavate a partire dal 1876 da A.~Goiran. Scavi successivi di De Stefani, De Mortillet, Tedeschi e Battaglia hanno via via sconvolto i depositi originali, rendendo difficile la ricostruzione stratigrafica. Il catalogo segnala anche una frequentazione nel Paleolitico superiore, che connette questa scheda al contesto più ampio della Valle dei Progni. Vedi anche la sezione Eneolitico~(\ref{app:eneolitico}).

\textbf{Negrar (VR).} \label{app:negrar} Abitato tra la fine del Neolitico medio e l'Eneolitico, con fasi culturali riferibili ai Vasi a Bocca Quadrata, alla cultura ``Fontbouisse'', allo stile campaniforme e allo ``White Ware''. La varietà degli aspetti culturali indica un luogo di incontro tra influenze diverse, coerente con la posizione della Valpolicella come zona di passaggio tra la pianura padana e le Alpi. Vedi anche la sezione Eneolitico~(\ref{app:eneolitico}).

\textbf{Rivoli Veronese (VR).} \label{app:rivoli} La Rocca di Rivoli, che domina la stretta dell'Adige, ha conservato tracce di insediamento e necropoli neolitici (Vasi a Bocca Quadrata e Lagozza), seguiti da occupazioni dell'antico e tardo Bronzo (cultura di Polada, fase Protovillanoviana) e poi da un abitato altomedievale e un castello bassomedievale. Le ricerche di G.~Pellegrini e L.H.~Barfield hanno identificato sei aree distinte sul sito. Vedi anche la sezione Bronzo~(\ref{app:rivoli_bronzo}).

\textbf{Sant'Anna d'Alfaedo (VR).} \label{app:alfaedo_neo} Vedi la descrizione completa alla sezione Paleolitico. La fase neolitica (tecnica Campignana) si innesta su una lunga sequenza che sale dal Paleolitico inferiore. Le attività di scheggiatura della selce continuano sino all'età del rame, rendendo questo sito uno dei più rappresentativi per la comprensione dell'economia litica della Lessinia.

\textbf{Sant'Ambrogio di Valpolicella (VR).} \label{app:santambrogio_neo} Sito pluristratificato con fase della cultura dei Vasi a Bocca Quadrata, cultura retica e periodo romano. La scheda è sintetica; la continuità dalla preistoria all'età romana è comunque significativa per comprendere la longevità dell'insediamento in Valpolicella.

\subsection{Pianura Veronese e Bassa}

\textbf{Erbè (VR).} \label{app:erbe} Sito pluristratificato con fasi neolitiche (Vasi a Bocca Quadrata), Bronzo antico e paleoveneto. La scheda è sommaria; il sito si colloca nella pianura irrigua veronese.

\textbf{Lavagno (VR).} \label{app:lavagno} Sito pluristratificato che copre un arco cronologico eccezionalmente lungo: le evidenze registrate partono dall'Eneolitico e si estendono attraverso il Bronzo tardo, l'età del Ferro, l'epoca romana, il Medioevo e il periodo austriaco. La collocazione in pianura veronese, in una zona di bassa collina, ha favorito la sedimentazione di cultura su cultura. Vedi anche le sezioni Eneolitico~(\ref{app:eneolitico}), Bronzo~(\ref{app:lavagno_bronzo}) e Ferro~(\ref{app:lavagno_ferro}).

\subsection{Prealpi Trevigiane e Pedemontana}

\textbf{Cornuda (TV).} Insediamento neolitico recente (cultura dei Vasi a Bocca Quadrata, fase più tarda con influenze Chassey-Lagozza) identificato dopo una frana che ha rimesso in luce materiali in giacitura secondaria. Il recupero è avvenuto in parte durante la rimozione di terre accumulate durante la seconda guerra mondiale. Alcune evidenze indicano una persistenza del sito nel Neolitico finale.

\textbf{Revine Lago (TV).} \label{app:revine} Villaggio su bonifica di tardo Neolitico e inizio Eneolitico (fine del quarto-inizi del terzo millennio a.p.z.a.) sul bordo dei Laghi di Revine, con frequentazione più sporadica nel corso dell'antico Bronzo (attorno al diciottesimo secolo a.p.z.a.). I primi materiali furono recuperati nel 1923 durante scavi per un canale; successivi scavi clandestini nel 1987 spinsero la Soprintendenza a organizzare campagne sistematiche nel 1992 e 1996. Il sito ha restituito materiali fittili, litici, resti faunistici e paleobotanici di grande interesse per la ricostruzione dell'economia e dell'ambiente preistorici. Dal fondo lacustre provengono anche armi in bronzo della media età del Bronzo (spade tipo Sauerbrunn, pugnale tipo Peschiera), che indicano una persistenza o una rifrequentazione del luogo. Vedi anche le sezioni Eneolitico~(\ref{app:eneolitico}) e Bronzo~(\ref{app:revine_bronzo}).

\subsection{Prealpi Vicentine e Piccole Dolomiti}

\textbf{Cornedo Vicentino (VI).} Riparo sotto roccia neolitico (ambito proto-veneto) probabilmente attribuibile a un abitato, senza strutture identificate. Il sito è uno dei pochi segnalati nelle Prealpi vicentine per questo periodo.

\subsection*{Nota di sintesi -- Neolitico}

Con ventidue schede il Neolitico è il secondo periodo per numerosità nel catalogo veneto (dopo il Bronzo). La distribuzione mostra una forte concentrazione nella pianura e nella bassa collina (Berici, bassa Lessinia, pianura veronese), che rispecchia sia la vocazione agricola delle comunità neolitiche sia la maggiore accessibilità di quelle aree per i ricercatori. I siti perilacustri (Fimon, Revine Lago) e quelli in ambiente umido (Este) indicano una preferenza per le risorse idriche. Le quote montane (Monte Avena, Mandriz, Cornedo Vicentino) suggeriscono che il territorio alpino fosse già percorso da pastori e cacciatori neolitici, ma le schede sono sparse e la ricerca di alta quota è quasi assente. Il grande complesso di Fimon, con le sue tre schede distinte, è di fatto il sito neolitico meglio documentato dell'intera regione.

% ============================================================
\section{Eneolitico}
\label{app:eneolitico}
% 9 siti (escludendo i rimandi)
% ============================================================

L'Eneolitico (età del Rame, circa 3.500-2.300 a.p.z.a.) è il periodo di diffusione della metallurgia del rame e dei primi metalli, con continuità nelle tradizioni neolitiche ma nuove reti di scambio a lunga distanza. Nel Veneto è spesso difficile distinguere Neolitico finale ed Eneolitico, e molte schede li trattano insieme.

\subsection{Bellunese e Feltrino}

\textbf{Fonzaso (BL) e Pedavena (BL) -- Monte Avena.} Vedi la descrizione completa nella sezione Neolitico~(\ref{app:avena_fonzaso}). La fase eneolitica di Monte Avena è interpretata come legata sia all'estrazione e lavorazione della selce sia alla caccia, in uno dei pochi contesti di alta quota catalogati per questo periodo nel Veneto.

\subsection{Cadore e Comelico}

\textbf{Selva di Cadore (BL) -- Riparo Mandriz.} Vedi la descrizione completa nella sezione Neolitico~(\ref{app:mandriz}). La frequentazione eneolitica si affianca a quella tardo-neolitica senza soluzione di continuità, indicando un utilizzo pluridecennale del riparo come base per pastori transumanti.

\subsection{Colli Euganei e Berici}

\textbf{Este (PD).} Vedi la descrizione nella sezione Neolitico~(\ref{app:este}). La fase eneolitica rientra nella lunga continuità insediativa del sito presso lo scolo delle Monache.

\subsection{Lago di Garda e Adige}

\textbf{Garda (VR).} Vedi la descrizione nella sezione Neolitico~(\ref{app:garda_neo}). La fase eneolitica fa parte della lunga stratificazione del promontorio gardesano.

\subsection{Lessinia}

\textbf{Fumane (VR).} Vedi la descrizione nella sezione Neolitico~(\ref{app:fumane_neo}). La fase eneolitica del sito di Fumane (seconda scheda) si affianca a quella neolitica.

\textbf{Negrar (VR).} Vedi la descrizione nella sezione Neolitico~(\ref{app:negrar}). L'abitato tra Neolitico recente ed Eneolitico è il nucleo principale del sito.

\subsection{Pianura Veronese e Bassa}

\textbf{Lavagno (VR).} Vedi la descrizione nella sezione Neolitico~(\ref{app:lavagno}). L'Eneolitico è la fase iniziale registrata nel catalogo per questo sito pluristratificato.

\subsection{Prealpi Trevigiane e Pedemontana}

\textbf{Revine Lago (TV).} Vedi la descrizione nella sezione Neolitico~(\ref{app:revine}). Il tardo Neolitico e l'Eneolitico costituiscono la fase principale del villaggio lacustre di Colmaggiore.

\subsection*{Nota di sintesi -- Eneolitico}

L'Eneolitico veneto nel catalogo MiC è in larga parte documentato come prolungamento o fase terminale dei siti neolitici già noti: quasi tutte le schede eneolitiche sono in realtà schede pluristratificate con fasi precedenti o successive. Il dato più interessante ai fini di questo libro è la presenza di siti in quota (Monte Avena, Riparo Mandriz) che indicano una mobilità verticale già strutturata durante l'età del Rame, con una frequentazione stagionale delle aree montane non dissimile da quella che potremmo aspettarci per i contesti cultuali alpini. La totale assenza di arte rupestre eneolitica catalogata in Veneto, a fronte della sua abbondanza in contesti analoghi (Valcamonica, Monte Bego), costituisce uno dei paradossi centrali di questo libro.

% ============================================================
\section{Età del Bronzo}
\label{app:bronzo}
% 42 siti
% ============================================================

L'età del Bronzo (circa 2.300-900 a.p.z.a.) è il periodo meglio rappresentato nel catalogo veneto, con quarantadue schede. È l'epoca dei castellieri, delle palafitte gardesane, degli insediamenti arginati di pianura e dei primi sistemi territoriali complessi. I metalli circolano su lunga distanza; l'allevamento transumante porta i pastori verso le quote alpine.

\subsection{Bellunese e Feltrino}

\textbf{Sedico (BL) -- Noàl.} \label{app:sedico} Castelliere molto attivo nell'età del Bronzo recente, inserito in una rete di siti che controllavano i versanti della valle del Piave. Attorno al nono-ottavo secolo a.p.z.a. Noàl faceva parte di un sistema di altura coordinato per il controllo delle rotte metallurgiche. La fase medievale (castello con torre) sfruttò la stessa posizione strategica, a controllo dell'imboccatura della valle del Cordevole.

\textbf{Limana (BL) -- San Pietro in Tuba.} \label{app:limana} Sito pluristratificato che copre la transizione Bronzo Recente-Finale e la prima età del Ferro, poi rioccupato in pieno Medioevo come castello di frontiera tra Trevigiano e Bellunese (distrutto nel 1366). Il dato più rilevante per il tema del libro è la segnalazione di un possibile contesto di deposizioni cultuali di manufatti in bronzo nei pressi di una sorgente, un tipo di evidenza che ricorre in molti santuari protostorici alpini. Vedi anche la sezione Ferro~(\ref{app:limana_ferro}).

\textbf{San Gregorio nelle Alpi (BL) -- Castel de Pedena.} \label{app:pedena} Sito d'altura con una prima fase del Bronzo medio-recente senza strutture difensive permanenti, a cui seguono, attorno al dodicesimo-undicesimo secolo a.p.z.a., strutture difensive vere e proprie che si mantengono fino alla prima età del Ferro. Dopo un lungo iato, il colle conosce una rioccupazione medievale strategica. Vedi anche la sezione Ferro~(\ref{app:pedena_ferro}).

\textbf{Belluno -- Col del Buson.} \label{app:colbuson} Abitato d'altura di età tardo-neolitica ed eneolitica (la scheda compare sotto il Bronzo nel catalogo per le fasi del Bronzo finale e dell'età del Ferro), integrato nell'ambiente circostante con pascoli e terreni agricoli. La documentazione di una struttura per la lavorazione della selce e di indizi di metallurgia in situ indica una comunità autosufficiente anche tecnologicamente. Vedi anche la sezione Ferro~(\ref{app:colbuson_ferro}).

\textbf{Belluno -- Sant'Anna di Pedecastello.} Tracce di frequentazione del Bronzo medio-recente ai margini della sommità del colle, con strutture principali attribuibili all'alto Medioevo e al Medioevo pieno. Le evidenze bronzee sono marginali rispetto alle fasi successive.

\subsection{Cadore e Comelico}

\textbf{San Vito di Cadore (BL) -- Sito VF1.} Definito di ``straordinaria importanza'' per la presenza di una sepoltura castelnoviana (Mesolitico), il sito ha visto anche frequentazioni dell'età del Bronzo e in epoca sub-attuale. La sovrapposizione cronologica su un arco di quasi diecimila anni ne fa un punto di riferimento per la comprensione della mobilità umana in Cadore.

\subsection{Colli Euganei e Berici}

\textbf{Lozzo Atestino (PD).} Necropoli a incinerazione del Bronzo finale; la scoperta di almeno due sepolture e di chiazze di terreno scuro nei campi vicini suggerisce l'esistenza di un'area funeraria più estesa nella zona di Vignalon.

\textbf{Este (PD).} \label{app:este_bronzo} Vedi la descrizione nella sezione Neolitico~(\ref{app:este}). Le almeno tre fasi bronzee del sito mostrano una riorganizzazione progressiva degli spazi insediativi.

\textbf{Montegrotto Terme (PD) -- Insediamento.} Frequentazione della media-tarda età del Bronzo (diciassettesimo-dodicesimo secolo a.p.z.a.) su un'area abbastanza estesa, con probabili attività agro-pastorali.

\textbf{Montegrotto Terme (PD) -- Villa ex Piacentini.} \label{app:montegrotto} Sito pluristratificato con prima frequentazione sporadica nell'età del Rame (prima metà del terzo millennio a.p.z.a.), insediamento stabile del Bronzo (quattordicesimo-dodicesimo secolo a.p.z.a.), villa romana monumentale di prima-quarta età imperiale, piccolo villaggio con necropoli altomedievale (ottavo-nono secolo) e insediamento produttivo più stabile tra il nono e il quattordicesimo secolo. La villa romana, interpretata come \textit{villa d'otium} per un committente di altissimo rango, è l'elemento più eccezionale della sequenza.

\subsection{Lago di Garda e Adige}

\textbf{Lazise (VR) -- Sito 1.} \label{app:lazise1} Insediamento palafitticolo con due fasi: una più antica tra il 1939 e il 1844 a.p.z.a. (analisi dendrocronologiche), l'altra della media età del Bronzo, con una terza frequentazione nell'età del Ferro. I materiali sono in parte frutto di scavi clandestini; un saggio sistematico è stato avviato nel 1986. Vedi anche la sezione Ferro~(\ref{app:lazise_ferro}).

\textbf{Garda (VR) -- Petroglifi di Monte Luppia.} \label{app:garda_petroglifi} La scheda MiC relativa ai petroglifi è quella di un ``sito pluristratificato'' del Bronzo e del Ferro, con la sintetica interpretazione: ``Petroglifi databili genericamente tra l'età del bronzo e quella del ferro.'' Monte Luppia è il sito di arte rupestre più noto del Veneto orientale, con incisioni di tipo geometrico e figure stilizzate su roccia calcarea affiorante. La catalogazione generica riflette le difficoltà di datazione delle incisioni in assenza di contesti stratigrafici chiusi; la letteratura specialistica propende per una datazione prevalentemente al Bronzo, con possibili continuazioni nel Ferro. Vedi anche la sezione Ferro~(\ref{app:garda_ferro}).

\textbf{Lazise (VR) -- Sito 2.} Secondo insediamento palafitticolo a Lazise, distinto dal precedente, datato tra la fase recente dell'antico Bronzo e la media età del Bronzo, con un rilievo topografico di 88 pali emergenti dal fondale del lago di Garda.

\textbf{Peschiera del Garda (VR) -- Sito 1.} Insediamento palafitticolo tra l'antico e il Bronzo recente, parzialmente sconvolto da scavi clandestini.

\textbf{Peschiera del Garda (VR) -- Sito 2.} Insediamento palafitticolo di grande estensione con due fasi principali: antico Bronzo (analisi dendrocronologiche: circa 2060 a.p.z.a.) e media età del Bronzo (circa 1617 a.p.z.a.). Seicento campioni di legno recuperati.

\subsection{Lessinia}

\textbf{Fumane (VR) -- Insediamento Bronzo recente.} Abitato probabilmente fortificato del Bronzo recente, con muretti a secco forse di epoca preistorica. Sito distinto dalla grotta paleolitica e dall'insediamento neolitico-eneolitico sopra descritti.

\textbf{Negrar (VR) -- Castelliere.} Insediamento, probabile castelliere, del Bronzo medio (cultura di Polada). Scheda sintetica.

\textbf{Rivoli Veronese (VR).} \label{app:rivoli_bronzo} Vedi la descrizione nella sezione Neolitico~(\ref{app:rivoli}). Le fasi bronzee includono il Bronzo antico (cultura di Polada) e il Bronzo finale (cultura Protovillanoviana).

\textbf{Sant'Anna d'Alfaedo (VR) -- Insediamento Bronzo.} Abitato probabilmente fortificato del Bronzo medio e recente, distinto dal sito paleolitico-neolitico della stessa area.

\subsection{Pianura Veronese e Bassa}

\textbf{Monteforte d'Alpone (VR).} \label{app:monteforte} Insediamento d'altura con sequenza dal Bronzo recente-finale all'età del Ferro, con fasi abitative documentate al dodicesimo-decimo secolo a.p.z.a., al nono, all'ottavo-settimo, al sesto-quinto e al secondo-primo secolo a.p.z.a. La continuità di occupazione per oltre un millennio è significativa. Vedi anche la sezione Ferro~(\ref{app:monteforte_ferro}).

\textbf{Nogarole Rocca (VR).} Insediamento tra la media età del Bronzo e il Bronzo finale. Nel 1977 fu recuperata una figurina fittile zoomorfa del Bronzo recente; nel corso dei sopralluoghi degli anni Settanta anche un coltello in bronzo tipo Vadena del Bronzo finale avanzato.

\textbf{Bovolone (VR).} Abitato del Bronzo con struttura abitativa a capanna, datato alla fase iniziale della media età del Bronzo.

\textbf{Colognola ai Colli (VR).} \label{app:colognola} Abitato su altura con fase del Bronzo medio e poi presenza della seconda età del Ferro (cultura retica e paleoveneta), seguita da un insediamento medievale. Vedi anche la sezione Ferro~(\ref{app:colognola_ferro}).

\textbf{Erbè (VR).} Vedi la descrizione nella sezione Neolitico~(\ref{app:erbe}). La fase bronzea si innesta sull'insediamento neolitico e precede quella paleoveneta.

\textbf{Legnago (VR).} Sito pluristratificato con Bronzo recente e cultura paleoveneta. Scheda sintetica.

\textbf{Roncà (VR).} \label{app:ronca_bronzo} Vedi la descrizione nella sezione Paleolitico~(\ref{app:ronca_paleo}). Le fasi del Bronzo recente-finale si inseriscono nella lunga stratificazione del sito.

\textbf{Villa Bartolomea (VR) -- Sito 1.} Sito pluristratificato con fase del Bronzo recente e periodo romano.

\textbf{Villa Bartolomea (VR) -- Sito 2.} Sito pluristratificato con fase del Bronzo medio-recente e periodo romano.

\textbf{Cerea (VR).} Sito palafitticolo in pianura veronese. Scheda sintetica (``sito palafitticolo'').

\textbf{Caldiero (VR).} Abitato d'altura protostorico (Bronzo medio-finale e cultura paleoveneta) con successivo insediamento medievale e post-medievale.

\textbf{Lavagno (VR).} \label{app:lavagno_bronzo} Vedi la descrizione nella sezione Neolitico~(\ref{app:lavagno}). La fase del Bronzo tardo si inserisce nella lunga stratificazione del sito.

\textbf{Nogara (VR).} Vasta zona a rischio archeologico con fase del Bronzo medio-recente, romana e medievale. La scheda è una ``zona di interesse'' senza strutture puntualmente identificate.

\subsection{Pianura e Delta}

\textbf{Dolo (VE).} Tracce di frequentazione del Bronzo recente 1 (seconda metà del quattordicesimo-inizio tredicesimo secolo a.p.z.a.) emerse durante i lavori di scavo per il Passante di Mestre tra il 2011 e il 2012. Il sito mostra attività agricole e di allevamento (canalette, pozzi, pozzetti, uno scarico di ceramica, una tomba di infante), poi sigillato da un evento alluvionale. Una frequentazione agricola sporadica dell'età del Ferro è successiva all'abbandono bronzeo.

\textbf{Arquà Petrarca (PD).} Insediamento palafitticolo perilacustre della cultura di Polada (Bronzo antico), ai piedi dei Colli Euganei.

\textbf{San Martino di Lupari (PD).} Insediamento di abitato arginato dell'ambito culturale Sub-appenninico (Bronzo recente-finale). La struttura ad argine trova confronti con gli insediamenti della pianura friulana.

\subsection{Prealpi Trevigiane e Pedemontana}

\textbf{Castello di Godego (TV).} Abitato fortificato del Bronzo medio-recente a forma quadrilatera con argine (``Le Motte''), che si estende in parte nel comune di San Martino di Lupari. Interpretato come vallo romano o medievale fino agli anni Settanta del Novecento, fu riconosciuto come preistorico solo dagli scavi del 1984-1985.

\textbf{Cordignano (TV).} \label{app:cordignano_bronzo} Le testimonianze per le fasi più antiche (Bronzo avanzato sub-appenninico e Bronzo finale protovillanoviano) sono frammentarie. Il sito diventa pienamente leggibile solo con la grande area sacra della prima età del Ferro e dell'età del Ferro veneto-antica (sesto-quarto secolo a.p.z.a.), attiva fino al quarto secolo. Vedi anche la sezione Ferro~(\ref{app:cordignano_ferro}).

\textbf{Montebelluna (TV).} \label{app:montebelluna_bronzo} Sito pluristratificato con frequentazioni documentate dal Bronzo medio-recente all'età moderna. Le fasi più significative nel catalogo sono quelle del Ferro e romane (complessi abitativi, impianti artigianali, necropoli), ma le radici bronzee del sito sono documentate. Vedi anche la sezione Ferro~(\ref{app:montebelluna_ferro}).

\textbf{Revine Lago (TV).} \label{app:revine_bronzo} Vedi la descrizione nella sezione Neolitico~(\ref{app:revine}). Le armi in bronzo recuperate dal fondo lacustre (spade tipo Sauerbrunn, pugnale tipo Peschiera) indicano frequentazioni anche nella media età del Bronzo.

\textbf{San Zenone degli Ezzelini (TV).} Insediamento d'altura a frequentazione stagionale di gruppi agro-pastorali, datato alla media età del Bronzo avanzata e al Bronzo medio-recente (quattordicesimo-tredicesimo secolo a.p.z.a.). I materiali provengono dalla sommità del colle e dalle pendici.

\textbf{Tarzo (TV).} \label{app:tarzo_bronzo} Vedi la descrizione nella sezione Paleolitico~(\ref{app:tarzo_paleo}). La fase bronzea aggiunge una terza frequentazione documentata per questo sito pluristratificato.

\textbf{Vedelago (TV).} Tracce di frequentazione di gruppi umani della cultura sub-appenninica nel Bronzo recente (quattordicesimo-tredicesimo secolo a.p.z.a.).

\subsection*{Nota di sintesi -- Età del Bronzo}

Con quarantadue schede, il Bronzo è il periodo più documentato nel catalogo veneto. La distribuzione geografica è molto più equilibrata rispetto alle epoche precedenti: palafitte sul Garda, castellieri nella pianura veronese e nelle Prealpi, insediamenti arginati nella pianura trevigiana, siti d'altura nel Bellunese. La sistematica presenza di controllo sui passi e sulle valli fluviali (Noàl sulla val Cordevole, Castel de Pedena, Rivoli sull'Adige, Monteforte sull'Alpone) indica un territorio organizzato in reti di sorveglianza e scambio. Il dato più rilevante per il tema di questo libro è la quasi totale assenza di arte rupestre catalogata per questo periodo nel Veneto, in contrasto con la ricchezza delle incisioni camune e liguri coeve: i petroglifi di Garda (Monte Luppia) sono l'unica eccezione esplicita nel catalogo, e anche lì la scheda è sintetica e la datazione ``generica''.

% ============================================================
\section{Età del Ferro}
\label{app:ferro}
% 31 siti
% ============================================================

L'età del Ferro (circa 900-0 a.p.z.a.) coincide nel Veneto con lo sviluppo della cultura paleoveneta (Este) e retica, con i suoi santuari delle acque, le necropoli a incinerazione e la metallurgia del ferro. È il periodo che precede la romanizzazione e che produce, nel catalogo MiC, la maggior varietà di tipologie di siti: insediamenti, necropoli e, notevolmente, santuari con depositi votivi.

\subsection{Bellunese e Feltrino}

\textbf{Mel (BL) -- Necropoli.} \label{app:mel_nec} Necropoli dell'età del Ferro a incinerazione con tombe a cassetta litica sormontate da tumuli sepolcrali, simile tipologicamente alle necropoli coeve di Este. Testimonianza dell'area di influenza paleoveneta nel Bellunese.

\textbf{Limana (BL).} \label{app:limana_ferro} Vedi la descrizione nella sezione Bronzo~(\ref{app:limana}). La fase del Ferro continua la vita del castelliere prima dell'abbandono e della successiva rioccupazione medievale. Il possibile deposito cultuale presso la sorgente rimane il dato più suggestivo del sito.

\textbf{San Gregorio nelle Alpi (BL).} \label{app:pedena_ferro} Vedi la descrizione nella sezione Bronzo~(\ref{app:pedena}). Il castelliere di Castel de Pedena è attivo fino alla prima età del Ferro.

\textbf{Belluno -- Col del Buson.} \label{app:colbuson_ferro} Vedi la descrizione nella sezione Bronzo~(\ref{app:colbuson}). La fase del Ferro completa la sequenza dell'abitato d'altura bellunese.

\textbf{Sedico (BL) -- Santuario di Libano, loc. Costiet.} Struttura in pietra a secco con materiali metallici, tra cui laminette in bronzo tipiche dei contesti cultuali venetici (tipo Lagole), una moneta romana in bronzo e oggetti in ferro. L'interpretazione è quella di un piccolo santuario frequentato durante l'età del Ferro e l'età romana. La mancanza di elementi stratigrafici chiari impedisce di stabilire se la frequentazione fosse continua o con interruzioni. La tipologia delle laminette collega questo sito alla rete di santuari delle acque venetici del Veneto orientale.

\textbf{Mel (BL) -- Abitato.} Abitato della seconda età del Ferro, con alcune strutture interpretabili come unità abitative (proprietà Calzavara e Rui-Sciulli) e numerose strutture lapidee probabilmente legate al drenaggio del substrato impermeabile. La documentazione disponibile consente di riconoscere un areale di diffusione antropica nell'attuale località Cioppa.

\subsection{Cadore e Comelico}

\textbf{Lozzo di Cadore (BL) -- Riva del Brodevin.} Necropoli con due fasi d'uso: la prima tra il settimo e il quarto secolo a.p.z.a., la seconda tra la fine del primo secolo a.p.z.a. e la seconda metà del quarto secolo. La continuità d'uso tra il Ferro e l'età tardo-romana indica la persistenza dell'abitato sottostante.

\textbf{Calalzo di Cadore (BL) -- Santuario di Lagole.} \label{app:lagole} Santuario territoriale venetico con continuità di frequentazione dalla metà del sesto secolo a.p.z.a. fino al quarto secolo. Il catalogo segnala una componente celtica nelle prime fasi e un culto legato alle acque curative, testimoniato da numerosi ex-voto anatomici. L'assenza di figurine femminili suggerisce che il culto fosse prevalentemente rivolto alla componente maschile della popolazione. Lagole è il sito rituale più importante del Cadore: la sua esistenza, insieme al santuario di Costiet (Sedico), indica che anche le vallate alpine erano attraversate da reti di sacralità riconoscibile.

\subsection{Colli Euganei e Berici}

\textbf{Barbarano Vicentino (VI).} Terrazzamenti artificiali attribuiti all'ambito veneto pre-protostorico. Scheda sintetica senza ulteriori dettagli.

\textbf{Montegrotto Terme (PD) -- Santuario.} Area ad ovest del Colle di S.~Pietro Montagnone con luogo di culto delle acque curative, attivo tra il settimo e il terzo secolo a.p.z.a. e testimoniato dagli ex-voto anatomici. Il collegamento tra sacralità delle acque termali e culto protostorico è significativo in un'area che diventerà poi famosa in età romana per i bagni di Aquae Patavinae.

\subsection{Lago di Garda e Adige}

\textbf{Lazise (VR).} \label{app:lazise_ferro} Vedi la descrizione nella sezione Bronzo~(\ref{app:lazise1}). La frequentazione del Ferro si aggiunge alle fasi palafitticole bronzee.

\textbf{Garda (VR) -- Petroglifi di Monte Luppia.} \label{app:garda_ferro} Vedi la descrizione nella sezione Bronzo~(\ref{app:garda_petroglifi}). I petroglifi sono databili genericamente tra l'età del Bronzo e quella del Ferro.

\textbf{Garda (VR) -- Sito pluristratificato.} Vedi la descrizione nella sezione Neolitico~(\ref{app:garda_neo}). La fase del Ferro (Veneti antichi e seconda età del Ferro) fa parte della lunga stratificazione del promontorio.

\subsection{Lessinia}

\textbf{Velo Veronese (VR).} Abitato d'altura probabilmente circondato da un vallo, dell'età del Ferro seconda fase. Due asce ritrovate nel 1854 durante la fondazione della chiesa locale; ceramiche eneolitiche e ossa recuperati nel 1948 sulla sommità; deposito di selci raccolto quasi sulla vetta del monte.

\subsection{Pianura Veronese e Bassa}

\textbf{Monteforte d'Alpone (VR).} \label{app:monteforte_ferro} Vedi la descrizione nella sezione Bronzo~(\ref{app:monteforte}). Le fasi del Ferro coprono l'ottavo-settimo e il sesto-quinto secolo a.p.z.a., fino alla fine dell'età del Ferro al secondo-primo secolo a.p.z.a.

\textbf{Cologna Veneta (VR).} Abitato e vasta necropoli paleoveneta; la scheda segnala anche una fase medievale. Sito rilevante per la rete funeraria del territorio paleoveneto nella pianura orientale veronese.

\textbf{Colognola ai Colli (VR).} \label{app:colognola_ferro} Vedi la descrizione nella sezione Bronzo~(\ref{app:colognola}). La seconda età del Ferro (cultura retica e paleoveneta) è la fase meglio documentata dell'abitato d'altura.

\textbf{Gazzo Veronese (VR) -- Abitato.} Abitato dell'età del Ferro. Scheda sintetica.

\textbf{Gazzo Veronese (VR) -- Necropoli.} Necropoli paleoveneta del periodo Atestino III, nella pianura veronese meridionale.

\textbf{Roncà (VR).} \label{app:ronca_ferro} Vedi la descrizione nella sezione Paleolitico~(\ref{app:ronca_paleo}). La tarda età del Ferro e la fase di romanizzazione completano la lunga sequenza del sito.

\textbf{Verona (VR).} Area pluristratificata nel centro storico di Verona, con fasi dall'età del Ferro (paleoveneto e retico) all'epoca moderna. La scheda catalogata testimonia la profondità stratigrafica del centro urbano.

\textbf{Lavagno (VR).} \label{app:lavagno_ferro} Vedi la descrizione nella sezione Neolitico~(\ref{app:lavagno}). La fase del Ferro si inserisce nella lunga stratificazione del sito pianeggiante.

\subsection{Pianura e Delta}

\textbf{Montagnana (PD).} Necropoli e abitato preromani (ambito veneto protostorico), con possibile villa rustica romana. La scheda segnala un'area funeraria del Ferro integrata nel paesaggio di pianura.

\textbf{Teolo (PD).} Materiali sporadici dell'ambito veneto protostorico, interpretati come tracce di frequentazione e/o di abitato. Scheda minima.

\subsection{Prealpi Trevigiane e Pedemontana}

\textbf{Cordignano (TV).} \label{app:cordignano_ferro} Vedi la descrizione nella sezione Bronzo~(\ref{app:cordignano_bronzo}). Il santuario territoriale di Cordignano diventa pienamente operativo a partire dalla fine del sesto secolo a.p.z.a. e rimane attivo fino al quarto secolo, quando gli editti di Teodosio vietano i culti pagani. Gli scavi dell'Università di Padova (1997-2006) hanno definito con precisione la grande area sacra, caratterizzata da una continuità devozionale di circa mille anni. Si tratta probabilmente di un santuario di frontiera, dedicato a divinità non ancora identificate, con componenti culturali veneto-antiche.

\textbf{Montebelluna (TV).} \label{app:montebelluna_ferro} Vedi la descrizione nella sezione Bronzo~(\ref{app:montebelluna_bronzo}). Il sito pluristratificato di Montebelluna mostra le sue fasi più ricche e meglio documentate nell'età del Ferro e in età romana (complessi abitativi, impianti artigianali, necropoli), con una continuità che va dal Bronzo recente fino all'epoca moderna.

\subsection{Prealpi Vicentine e Piccole Dolomiti}

\textbf{Rotzo (VI).} Sito d'altura di ambito veneto-retico, con posizione strategica all'imboccatura della valle dell'Astico. Scheda sintetica.

\textbf{Creazzo (VI).} Sito pluristratificato di ambito veneto-romano, collegato alla coltivazione dei campi e alla lavorazione e stoccaggio dei prodotti agricoli. Scheda sintetica.

\textbf{Montecchio Maggiore (VI).} Area di frequentazione preromana (ambito veneto) e insediamento rustico romano, con necropoli longobarda. Continuità insediativa dal Ferro all'alto Medioevo.

\textbf{Santorso (VI).} Sito pluristratificato di ambito veneto-romano, con strutture dell'abitato preromano e romano. Scheda sintetica.

\textbf{Isola Vicentina (VI).} Sito pluristratificato di ambito veneto-romano. I resti romani includono un insediamento rustico con un settore interpretato come \textit{fullonica} (impianto per la lavorazione della lana). Scheda sintetica.

\subsection*{Nota di sintesi -- Età del Ferro}

L'età del Ferro nel catalogo veneto rivela la mappa di una società complessa, organizzata in reti di insediamenti e santuari. I dati più preziosi per il tema di questo libro sono i santuari delle acque: Montegrotto Terme (Colli Euganei), Lagole di Calalzo (Cadore), Costiet di Sedico (Bellunese), il possibile deposito votivo di Limana. Questi siti mostrano che la sacralizzazione di punti fisici nel paesaggio, in particolare quelli legati all'acqua, era una pratica strutturata e geograficamente distribuita nell'intero territorio veneto, dalle pianure termali alle vallate alpine. L'assenza di santuari analoghi nell'arco prealpino centrale e nelle Piccole Dolomiti non significa che non esistessero: significa che non li abbiamo trovati, o che non li abbiamo riconosciuti. Il santuario di Cordignano (al confine con il Friuli), attivo per circa mille anni, dimostra la capacità del registro archeologico di conservare luoghi di culto quando le condizioni topografiche e stratigrafiche sono favorevoli.

% ============================================================
\section*{Conclusione dell'appendice}
% ============================================================

I 136 siti censiti in queste pagine non sono un'enciclopedia della preistoria veneta, ma sono l'inventario di ciò che il catalogo ufficiale riconosce come noto e formalmente catalogato. Le lacune sono strutturali: mancano quasi completamente i siti d'alta quota (sopra i 1.500~m), le aree di transizione Bellunese-Cadore non sono sistematicamente prospettate, e intere epoche (Mesolitico in particolare) sono quasi prive di schede pur essendo rappresentate nei siti limitrofi del Trentino e del Friuli.

L'utilità di questa raccolta, come strumento per il ragionamento proposto nel libro, è duplice. Da un lato fornisce una base dati concreta che permette di distinguere il silenzio dell'arte rupestre veneta dal silenzio dell'archeologia in generale: il Veneto non è vuoto di preistoria, anzi è ricco di siti, di culture, di occupazioni ripetute; è specificamente l'arte rupestre a mancare, e questo la dice lunga sui tipi di ricerca che sono stati condotti e sui tipi che mancano. Dall'altro lato, i santuari delle acque dell'età del Ferro (Lagole, Montegrotto, Costiet, Limana) e i siti d'altura del Bronzo (Monte Avena, San Zenone, Noàl) indicano che il territorio montano veneto era percorso, frequentato e sacralizzato: le condizioni per la produzione di arte rupestre esistevano. Aspettano solo di essere cercate.
